\scp{Final vowel-glide sequences}{06-03}
The reflexes of final vowel-glide sequences *-aw
and *-ay in the Aru languages are diverse.
 In this section we only focus on the reflexes of final vowel-glide
sequences for words in which a medial consonant was preserved,
as the reflexes are clearest in this environment.
Reflexes of vowel-glide sequences include: preservation of the
full sequence (*-aw = \it{aw;} *-ay = \it{ay}),
fusion to a mid vowel (*-ay > \it{ɛ}),
deletion of the final glide (*-aw \it{>} -\it{a}),
and preservation of the final glide as a vowel (*-aw > \it{u}; *-ay > \it{i}).
The full range of reflexes found throughout the Wallacea region
are seen in different \tsc{Aru} languages.
No changes to vowel-glide sequences align with different nodes
within \tsc{Aru} and a single language does not necessarily have parallel changes
for *-aw and *-ay.
Examples of the reflexes of *-aw are given in \trf{06-16},
and *-ay in \trf{06-17}.

There is not enough data to determine a consistent reflex of
*-aw for Lorang, Barakai, Karey, Mariri, and Manombai.
There is also not enough data to determine a consistent reflex
of *-ay for Kompane,
Lola, Lorang, Koba, Barakai,
Karey, Mariri, and East Tarangan.
(Additionally, the reflexes of *bəRsay ‘paddle’ (\trf{06-13})
appear to be irregular with respect to the final vowel-glide
sequence.) 

\begin{table}\tcp{\tsc{Aru} *-aw = **-aw}{06-16}
\begin{adjustbox}{width=\textwidth}
	\begin{threeparttable}\st{2pt}
	\begin{tabular}{lllllll}\lsptoprule
local gl.	&		&	sago	&	\hp{*-}walk	&	{\hr}sun	&	\hp{*-}steal	&	house\\
PMP	&	\hr*-aw	&		&	\hp{*-}*lak\tbr{aw}	&	\hr*qaləj\tbr{aw}	&	\hp{*-}*nak\tbr{aw}	&	\\
Proto-Aru	&	**-aw	&	**sak\tbr{aw}	&	**-lak\tbr{aw}	&	**laR\tbr{aw}	&	**-naŋg\tbr{aw}	&	**pal\tbr{aw}\\ \midrule
Ujir	&	{\hr\hr}-aw, -u	&	{\hr\hr}sak\tbr{aw}	&		&	{\hr\hr}law\tbr{u}	&	{\hr\hr}-y-naŋ\tbr{aw}	&	\\
Kompane	&	{\hr\hr}-aw	&	{\hr\hr}ak\tbr{aw}	&		&	{\hr\hr}lah\tbr{aw}	&		&	\\
East Kola	&	{\hr\hr}-aw	&	{\hr\hr}ak\tbr{aw}	&	{\hr\hr}-lak\tbr{áu}	&	{\hr\hr}lah\tbr{aw}	&	{\hr\hr}-fanag\tbr{aw}	&	{\hr\hr}ɸal\tbr{aw}\\
Lola	&	{\hr\hr}-u , -a	&	{\hr\hr}yaʔ\tbr{u}	&		&	{\hr\hr}lah\tbr{u}	&		&	{\hr\hr}wal\tbr{a}sín\\
Lorang	&	{\hr\hr}(-a)	&	{\hr\hr}yak\tbr{a}	&		&		&		&	\\
Dobel	&	{\hr\hr}-u	&		&		&	{\hr\hr}lar\tbr{u}\su{‡}	&	{\hr\hr}-naŋ\tbr{u}	&	{\hr\hr}ɸal\tbr{u}\su{Δ}\\
Koba	&	{\hr\hr}-a	&	{\hr\hr}yak\tbr{a}	&		&	{\hr\hr}lar\tbr{a}	&		&	\\
Barakai	&	{\hr\hr}(-a)	&	{\hr\hr}ag\tbr{a}	&		&		&		&	\\
Karey	&	{\hr\hr}(-aw)	&	{\hr\hr}ag\tbr{aw}	&		&		&		&	\\
Batuley	&	{\hr\hr}-a	&	{\hr\hr}ag\tbr{a}	&		&	{\hr\hr}taɸer lar\tbr{a}	&		&	\\
Mariri	&	{\hr\hr}(-ao)	&	{\hr\hr}ág\tbr{ao}	&		&		&		&	\\
Manombai	&	{\hr\hr}(-a)	&	{\hr\hr}ek\tbr{a}	&		&		&		&	\\
NE. Tar.\su{†}	&	{\hr\hr}-a	&	{\hr\hr}ak\tbr{a}	&	{\hr\hr}-r-lɔk\tbr{a}	&		&	{\hr\hr}-naŋ\tbr{a}	&	{\hr\hr}hál\tbr{a}\\
SE. Tar.	&	{\hr\hr}-aw	&	{\hr\hr}ag\tbr{aw}	&	{\hr\hr}-r-lag\tbr{aw}	&		&		&	{\hr\hr}ɸal\tbr{aw}\\
NW. Tar. 1	&	{\hr\hr}-a	&	{\hr\hr}ak\tbr{a}	&	{\hr\hr}-r-lɔk\tbr{a}	&		&	{\hr\hr}-nɔŋ\tbr{a}	&	\\
NW. Tar. 2	&	{\hr\hr}-a	&	{\hr\hr}ak\tbr{a}	&	{\hr\hr}-r-lɔk\tbr{a}	&		&	{\hr\hr}-nɔŋ\tbr{a}	&	\\
SW. Tar. 1	&	{\hr\hr}-aw	&	{\hr\hr}ɔk\tbr{aw}	&	{\hr\hr}-r-lɔk\tbr{aw}	&		&	{\hr\hr}-nɔŋ\tbr{aw}	&	{\hr\hr}ɸal{\tl}ɸal\tbr{aw}\su{ǁ}\\
SW. Tar. 2	&	{\hr\hr}-aw	&	{\hr\hr}ɔk\tbr{aw}	&		&		&	{\hr\hr}-nɔŋ\tbr{aw}	&	\\
		\lspbottomrule
	\end{tabular}
		\begin{tablenotes}\tnsize
			\item[†]	{Data from varieties of \tsc{Tarangan}
								are as follows: NE. Tarangan, Jelia village;
								SE. Tarangan from Dosi Mar village;
								NW. Tarangan 1 from Doka Timur village;
								NW. Tarangan 2 from Lutur village;
								SW. Tarangan 1 from Ngaibor village; and
								SW. Tarangan 2 from Kalar-kalar village.}
			\item[‡]	{Lola \it{lahu} ‘sun’ is from Warabal village.
								Lola village has \it{laɸo} ‘sun’.}
			\item[Δ]	{Dobel \it{ɸalu} is ‘very large clan house’.}
			\item[{ǁ}]	{Rick Nivens (pers. comm.
								April 2024) doubts that Ngaibor \it{ɸal{\tl}ɸalaw} is a valid form
								and proposes that it is probably due to contact with SE. Tarangan.}
		\end{tablenotes}
	\end{threeparttable}
\end{adjustbox}
\end{table}

\begin{table}\tcp{\tsc{Aru} *-ay = **-ay}{06-17}
\begin{adjustbox}{width=\textwidth}
	\begin{threeparttable}\st{2pt}
	\begin{tabular}{llllllll}\lsptoprule
local gl.	&		&	{\hr}dry in sun	&	{\hr}climb	&	bitter	&	open	&	boat	&	sailboat\\
PMP	&	\hr*-ay	&	\hr*waRi	&	\hr*sak\tbr{ay}	&		&		&		&	\\
Proto-Aru	&	**-ay	&	**-waR\tbr{ay}	&	**sak\tbr{ay}	&	**mak\tbr{ay}	&	**pak\tbr{ay}	&	**kal\tbr{ay}	&	**let\tbr{ay}\\	\midrule
Ujir	&	{\hr\hr}-i, -ay	&	{\hr\hr}-wo\tbr{y}	&		&		&	{\hr\hr}ɸak\tbr{i}	&	{\hr\hr}kal\tbr{i}	&	{\hr\hr}let\tbr{ay}\\
Kompane	&	{\hr\hr}(-i)	&	{\hr\hr}-gwah\tbr{i}	&		&		&		&		&	\\
East Kola	&	{\hr\hr}-i	&	{\hr\hr}-weh	&	{\hr\hr}-as\tbr{i}	&	{\hr\hr}mak\tbr{i}	&	{\hr\hr}-ɸak\tbr{i}	&	{\hr\hr}kal\tbr{i}	&	{\hr\hr}let\tbr{i}\\
Lola	&	{\hr\hr}(-i)	&	{\hr\hr}-wah\tbr{i}	&		&		&		&		&	\\
Lorang	&	{\hr\hr}(-a)	&	{\hr\hr}-ger\tbr{a}	&		&		&		&		&	\\
Dobel	&	{\hr\hr}-ay	&	{\hr\hr}-kwar\tbr{i}	&	{\hr\hr}s{\tl}saʔ\tbr{ay}	&	{\hr\hr}maʔ\tbr{ay}	&	{\hr\hr}-ɸaʔ\tbr{ay}	&	{\hr\hr}ʔal\tbr{ay}	&	{\hr\hr}let\tbr{ay}\\
Koba	&	{\hr\hr}(-a)	&	{\hr\hr}-ʍer\tbr{a}	&		&		&		&		&	\\
Batuley	&	{\hr\hr}-ey	&	{\hr\hr}wa\tbr{ye}r\su{‡}	&		&	{\hr\hr}mag\tbr{ey}	&		&	{\hr\hr}kál\tbr{ey}	&	{\hr\hr}let\tbr{ey}\\
Manombai	&	{\hr\hr}-ay	&	{\hr\hr}-gar\tbr{ay}	&		&		&	{\hr\hr}ɸak\tbr{ay}	&		&	\\
NE. Tar.\su{†}	&	{\hr\hr}(-ɛ)	&	{\hr\hr}-gwar\tbr{ɛ}	&		&		&		&		&	\\
SE. Tar.	&	{\hr\hr}(-ay)	&	{\hr\hr}-gwar\tbr{ay}	&		&		&		&		&	\\
NW. Tar. 1	&	{\hr\hr}-ɛ	&	{\hr\hr}-gar\tbr{ɛ}	&		&	{\hr\hr}mak\tbr{ɛ}	&	{\hr\hr}-ɸák\tbr{ɛ}	&		&	{\hr\hr}lɛt\tbr{ɛ}\\
NW. Tar. 2	&	{\hr\hr}-ay	&	{\hr\hr}-gar\tbr{ay}	&	{\hr\hr}-ak\tbr{ay}	&		&		&		&	{\hr\hr}lɛt\tbr{ay}\\
SW. Tar. 1	&	{\hr\hr}-ay	&	{\hr\hr}-gar\tbr{ay}	&	{\hr\hr}-ak\tbr{ay}	&	{\hr\hr}mak\tbr{ay}	&		&		&	{\hr\hr}lɛt\tbr{ay}\\
SW. Tar. 2	&	{\hr\hr}-ay	&	{\hr\hr}-gar\tbr{ay}	&	{\hr\hr}-ak\tbr{ay}	&	{\hr\hr}mak\tbr{ay}	&	{\hr\hr}-ɸak\tbr{ay}	&	{\hr\hr}kal\tbr{ay}\su{Δ}	&	{\hr\hr}lɛt\tbr{ay}\\
		\lspbottomrule
	\end{tabular}
		\begin{tablenotes}\tnsize
			\item[†]	{\tsc{Tarangan} data is from the same varieties
								and villages as given in \trf{06-16}.
								Northwest Tarangan \it{lɛtay} ‘sailboat’ is from Rebi village.}
			\item[‡]	{Batuley \it{wayer} ‘dry in sun’ has metathesis
								of the medial and final consonants from **warey.}
			\item[Δ]	{Southwest Tarangan has \it{kalin} ‘boats’ as a plural.
								Rick Nivens (pers. comm.
								March 2021) further points out that this is unusual since the
								plural suffix \it{-in} is usually only found on inalienably possessed nouns.}
		\end{tablenotes}
	\end{threeparttable}
\end{adjustbox}
\end{table}
